\documentclass[11pt,a4paper]{article}
\usepackage[margin=22mm]{geometry}
\usepackage[T1]{fontenc}
\usepackage{lmodern}
\usepackage{microtype}
\usepackage[hidelinks]{hyperref}
\usepackage{parskip}
\setlength{\parindent}{0pt}
\pagestyle{plain}

\newcommand{\sourceheading}[2]{%
  \section*{#1}
  \textbf{#2}\par\medskip
}

\begin{document}

\begin{center}
  {\Large\bfseries LN905 Reading into Writing Practice}\par
  \vspace{4pt}
  {\large Gender -- 4 August 2026}
\end{center}

\section*{Question}

\textbf{To what extent does evidence from Norway show that mandatory corporate-board gender quotas are an effective way to advance workplace gender equality?}

Use all three extracts to develop a qualified answer. The extracts below are exam-style adaptations of the cited peer-reviewed papers. Their wording has been condensed and paraphrased while preserving the studies' central claims, evidence and limits.

\sourceheading{Extract 1}{Adapted from Ahern and Dittmar (2012), ``The Changing of the Boards: The Impact on Firm Valuation of Mandated Female Board Representation''}

Norway announced in 2003 that public limited companies would have to allocate roughly 40\% of board seats to each gender. Women then held only about 9\% of directorships, so firms with all-male or nearly all-male boards had to make the largest changes. Ahern and Dittmar use this pre-quota variation to estimate how an externally imposed change in board composition affected firms.

The authors first examine stock returns around the announcement. Among the firms for which the required information was available, companies with no female directors before the law experienced an average abnormal return of about -3.5\% over the five-day event window, while firms that already had female directors experienced little change. In panel analyses covering the years surrounding implementation, a larger quota-induced change was also associated with a larger fall in Tobin's Q, a market-based measure of firm value.

Board characteristics changed at the same time. The newly constituted boards were younger and had less chief-executive experience. The firms most affected by the quota also increased leverage and acquisitions, and their operating performance weakened. The authors interpret this combination as evidence that restricting shareholders' choice of directors imposed costs: the law accelerated appointments before firms had developed a sufficiently large pool of experienced candidates.

This evidence concerns firm valuation and governance, not gender equality directly. It can support the claim that a quota may have economic costs if the event-study and panel estimates isolate the law correctly. It cannot by itself establish whether more representative boards improve fairness, alter women's wider career opportunities, or produce longer-run benefits that market prices did not immediately capture.

\textit{Source: Quarterly Journal of Economics, 127(1), 137--197. DOI: \href{https://doi.org/10.1093/qje/qjr049}{10.1093/qje/qjr049}.}

\newpage

\sourceheading{Extract 2}{Adapted from Eckbo, Nygaard and Thorburn (2022), ``Valuation Effects of Norway's Board Gender-Quota Law Revisited''}

Eckbo and colleagues directly revisit the claim that Norway's quota reduced firm value. They argue that policy-event studies face three linked problems: it can be unclear which day conveyed new information about the quota; firms' returns on the same day are correlated rather than independent; and firm characteristics or economy-wide events can be mistaken for effects of the law.

Their re-examination begins with the earlier five-day announcement window. The strongest negative return occurred after a news report that appeared to lower, rather than raise, the probability of a binding quota. A negative return following that report is therefore difficult to interpret as a cost of imposing the law. In addition, the original significance test treated firms exposed to the same policy news as though their returns were independent. Once contemporaneous correlation is taken into account, the study cannot reject a zero market reaction.

The authors then use portfolio methods across five quota-related news events, with roughly 122--133 firms available at each event, and examine longer-run risk-adjusted returns and operating measures. These analyses do not produce reliable evidence of a negative valuation effect. They conclude that the supply of qualified women at the time was sufficient to meet the quota without imposing the shareholder costs previously claimed.

This is a methodological rebuttal, not proof that quotas create broad equality. A statistically insignificant valuation effect means the available design does not establish a cost; it does not demonstrate that the effect is beneficial or exactly zero. The evidence comes from one country's corporate and labour-market setting and says little by itself about recruitment, promotion or pay below board level.

\textit{Source: Management Science, 68(6), 4112--4134. DOI: \href{https://doi.org/10.1287/mnsc.2021.4031}{10.1287/mnsc.2021.4031}. CC BY 4.0.}

\newpage

\sourceheading{Extract 3}{Adapted from Bertrand, Black, Jensen and Lleras-Muney (2019), ``Breaking the Glass Ceiling? The Effect of Board Quotas on Female Labour Market Outcomes in Norway''}

Bertrand and colleagues ask whether the quota's direct effect on board composition spread to women elsewhere in the corporate labour market. They link Norwegian registers on companies, boards, earnings, education and employment over 1998--2014. Because some firms changed legal form after the policy announcement, the authors analyse both firms still subject to the quota and an intent-to-treat group of 581 companies that were public limited companies in 2003.

The direct representational change was large: the median share of women on affected boards rose from 0\% in 2003 to 40\% in 2008. Women appointed after the reform were, on observable measures, better qualified than earlier female directors. The gender earnings gap within boards also narrowed. These results weaken the claim that firms could comply only by appointing unsuitable candidates and show that a legal rule can change access to a clearly defined set of senior positions.

The wider effects were much smaller. The authors find no robust improvement for the larger group of women employed by quota-bound firms, and no clear gains for highly qualified women whose profiles resembled directors' but who did not enter boardrooms. Female enrolment in business education did not rise. Some early-career outcomes improved for female business graduates, but similar changes among science graduates make it difficult to attribute those gains to the quota.

The study therefore distinguishes a successful targeted intervention from a general workplace-equality policy. Seven years after compulsory implementation, the quota clearly changed who occupied board seats but had little discernible impact beyond the women appointed. Longer-run spillovers may still emerge, and changes in company legal form complicate estimation, but the observed evidence does not justify treating board representation as a substitute for reforms to hiring, promotion and pay.

\textit{Source: Review of Economic Studies, 86(1), 191--239. DOI: \href{https://doi.org/10.1093/restud/rdy032}{10.1093/restud/rdy032}.}

\end{document}
